\fichatitulo{37 · FUNDAMENTO}{JMLR · 2020}{Exploring the Limits of Transfer Learning}
{\small Raffel et al. \textit{Exploring the Limits of Transfer Learning with a Unified Text-to-Text Transformer}. JMLR · 2020. \href{https://jmlr.org/papers/volume21/20-074/20-074.pdf}{Fonte consultada}.\par}

\campo{Problema e objetivo}
Como comparar estratégias de transferência em tarefas heterogêneas?

\campo{Principal contribuição · síntese interpretativa}
Unifica problemas de linguagem no formato texto-para-texto e apresenta T5 e o corpus C4.

\campo{Método / natureza da evidência}
Estudo comparativo de objetivos, arquiteturas, dados e estratégias de transferência; o recorte consultado caracteriza o desenho e o modelo.

\campo{Resultado ou argumento principal}
A formulação comum permite expressar diferentes tarefas com entradas e saídas textuais e comparar escolhas dentro de um mesmo enquadramento.

\campo{Excertos-chave · seleção literal}
\begin{itemize}
\excerto{1}{a text-to-text format}{Resumo.}
\excerto{2}{pre-training objectives, architectures, unlabeled data sets, transfer approaches}{Resumo.}
\end{itemize}

\campo{Limitações e análise crítica}
Análise da ficha: esta leitura focal não reproduz todas as ablações. A uniformidade de formato não elimina diferenças entre critérios de avaliação ou necessidades dos usuários.

\campo{Aproveitamento na dissertação · proposta}
Fundamentação: situar modelos codificador-decodificador e separar formato da tarefa de qualidade da resposta.

\registro{Fonte consultada nesta preparação: Resumo, introdução e seção 2.1; consulta parcial. Chave: \texttt{raffelEtAl2020t5}.}
